\subsection{PRODUCT OF TWO SETS}

THEOREM 1. The relation

\[
(\forall \mathrm{X}) (\forall \mathrm{Y}) (\exists \mathrm{Z}) (\forall z) ((z \in \mathrm{Z}) \iff (\exists x) (\exists y) (z = (x, y) \text {   and   } x \in \mathrm{X} \text {   and   } y \in \mathrm{Y})
\]

is true. In other words, for all X and all Y the relation ``z is an ordered pair and \(pr_{1}z \in X\) and \(pr_{2}z \in Y\) '' is collectivizing in z.

Let \(A_y\) denote the set of all objects of the form \((x, y)\) for \(x \in X\) (cf.~§1, no. 6, criterion C53). Let \(R\) be the relation \(z \in A_y\) , which is equivalent to \((\exists x)(z = (x, y)\) and \(x \in X)\) . It is clear that the relation

\[
(\forall y) (\exists \mathrm{A}) (\forall z) (\mathrm{R} \Longrightarrow (z \in \mathrm{A}))
\]

is true, by virtue of S5 (Chapter I, §4, no. 2). It then follows from S8 that the relation \((\exists y)(y \in \mathbf{Y}\) and \(\mathbf{R})\) is collectivizing in \(z\) . But this relation is equivalent to \((\exists x)(\exists y)(y \in \mathbf{Y}\) and \(x \in \mathbf{X}\) and \(z = (x, y)\) ); hence the result.

DEFINITION 1. Given two sets X and Y, the set

\[
\mathcal {E} _ {z} ((\exists x) (\exists y) (z = (x, y) \text {   and   } x \in X \text {   and   } y \in Y))
\]

is called the product of X and Y and is denoted by \(X \times Y\) .

The relation \(z \in X \times Y\) is thus equivalent to `` \(z\) is an ordered pair and \(\text{pr}_1z \in X\) and \(\text{pr}_2z \in Y\) ''. The sets \(X\) and \(Y\) are called the first and second factors of \(X \times Y\) .

PROPOSITION 1. If A', B' are non-empty sets, the relation A' × B' ⊂ A × B is equivalent to ``A' ⊂ A and B' ⊂ B''.

In the first place, the relation \(z \in A' \times B'\) is equivalent to `` \(z\) is an ordered pair and \(\mathrm{pr}_1z \in A'\) and \(\mathrm{pr}_2z \in B'\)''; therefore, without any restriction on \(A'\) and \(B'\) , the relation `` \(A' \subset A\) and \(B' \subset B\)'' implies

\[
\mathbf {A} ^ {\prime} \times \mathbf {B} ^ {\prime} \subset \mathbf {A} \times \mathbf {B}.
\]

Conversely, let us first show that if \(B' \neq \emptyset\) (without restriction on \(A'\) ), the relation \(A' \times B' \subset A \times B\) implies \(A' \subset A\) . Let x be an element of \(A'\) ; since \(B' \neq \emptyset\) , there is an object y which is an element of \(B'\) ; we have \((x, y) \in A' \times B'\) , hence \((x, y) \in A \times B\) , and consequently \(x \in A\) ; thus \(A' \subset A\) . Similarly, if \(A' \neq \emptyset\) , the relation \(A' \times B' \subset A \times B\) implies \(B' \subset B\) . Hence the result.

PROPOSITION 2. Let A and B be two sets. The relation \(\mathbf{A} \times \mathbf{B} = \varnothing\) is equivalent to `` \(\mathbf{A} = \varnothing\) or \(\mathbf{B} = \varnothing\) ''.

For the relation \(z \in A \times B\) implies \(\mathsf{pr}_1z \in A\) and \(\mathsf{pr}_2z \in B\) ; hence \(A \neq \emptyset\) and \(B \neq \emptyset\) . Conversely, the relation `` \(x \in A\) and \(y \in B\) '' implies \((x, y) \in A \times B\) and hence \(A \times B \neq \emptyset\) . In other words, the relation \(A \times B \neq \emptyset\) is equivalent to `` \(A \neq \emptyset\) and \(B \neq \emptyset\) ''; hence the result.

If A, B, C are sets, we write

\[
(\mathbf {A} \times \mathbf {B}) \times \mathbf {C} = \mathbf {A} \times \mathbf {B} \times \mathbf {C};
\]

an element \(((x, y), z)\) of A × B × C is written \((x, y, z)\) and is called a triple. Again, if A, B, C, D are sets, we write

\[
(\mathbf {A} \times \mathbf {B} \times \mathbf {C}) \times \mathbf {D} = \mathbf {A} \times \mathbf {B} \times \mathbf {C} \times \mathbf {D};
\]

and so on.

